\documentclass[11pt]{article}
\usepackage[a4paper,margin=25mm]{geometry}
\usepackage[T1]{fontenc}
\usepackage{lmodern}
\usepackage{booktabs}
\usepackage{array}
\usepackage{longtable}
\usepackage{hyperref}
\usepackage{xcolor}
\usepackage{listings}
\lstset{basicstyle=\ttfamily\small, columns=fullflexible, keepspaces=true,
  breaklines=true, frame=single, framerule=0.3pt, rulecolor=\color{gray}}
\newcommand{\code}[1]{\texttt{#1}}
\newcommand{\us}{\textmu s}

\title{\code{JScan}: reading the BGO temperature-scan runs in Julia}
\author{BGOTSCAN}
\date{2026-10-14}

\begin{document}
\maketitle

\section{Purpose}

\code{JScan} is the Julia side of the BGOTSCAN analysis: the passes over the
full data set. This document covers the reader (opening a run, getting its
parameters, reading blocks of waveforms), the layout of the files written by
the acquisition, which both the Julia and the Python code rely on, and the
pass (section~\ref{sec:pass}): one read of a run that writes, per event, the
corrected summed waveform, the charge of every channel, and the
single-photoelectron candidates of every channel, for the offline analysis in
Python.

\section{The run files}

Each run is one HDF5 file written by WebDump, the recording software of the
CAEN DT2745 digitizer (Scope firmware): 64 channels, 125~MHz sampling, 16-bit
ADC. The file name encodes the run conditions,
\begin{lstlisting}
<T>K_<V>V_th<thr>.h5      e.g.  185K_39d5V_th19.h5
\end{lstlisting}
with \code{T} the crystal temperature (thermometer T3) in kelvin, \code{V} the
SiPM bias in volts (\code{d} stands for the decimal point: \code{39d5V} is
39.5~V) and \code{thr} the trigger threshold in ADC counts.

\subsection{Datasets}

\begin{center}
\begin{tabular}{llll}
\toprule
dataset & shape (as written) & type & content \\
\midrule
\code{waveforms}    & (events, channels, samples) & uint16  & raw ADC counts \\
\code{channels}     & (channels,)                 & int16   & digitizer channel of each row \\
\code{timestamp\_ns} & (events,)                  & float64 & trigger time of each event, ns \\
\code{trigger\_id}  & (events,)                   & uint32  & board trigger counter \\
\bottomrule
\end{tabular}
\end{center}

The shape above is the C order of h5py and numpy. HDF5 stores arrays in that
order, and Julia, which is column-major, sees the same bytes with the axes
reversed: in \code{JScan} a block of waveforms is a
(samples, channels, events) array, so \code{w[:, c, e]} is the waveform of
channel \code{c} in event \code{e}, contiguous in memory. The datasets are
chunked by two events and compressed with LZF (see section~\ref{sec:lzf}). For
the runs of 2026-10-02: 5000 events, 64 channels, 3000 samples.

\code{trigger\_id} counts every trigger the board saw; gaps between the ids of
consecutive recorded events are triggers that were not recorded.

\subsection{Attributes}

All attributes are on the root group. Table~\ref{tab:attrs} lists them with
the value of the run \code{128K\_37V\_th19}. Arrays have one entry per
recorded channel.

\begin{longtable}{l l p{72mm}}
\caption{Root attributes of a run file.}\label{tab:attrs}\\
\toprule
attribute & example & meaning \\
\midrule
\endfirsthead
\toprule
attribute & example & meaning \\
\midrule
\endhead
\code{model}, \code{serial}  & DT2745, 52694 & digitizer \\
\code{fw\_type}, \code{fpga\_fw} & Scope, 1.0.257 & firmware \\
\code{software}             & WebDump & recording program \\
\code{url}                  & dig2://caendgtz-usb-52694 & connection \\
\code{start\_time}          & 2026-10-02T14:47:45 & start of the run \\
\code{events}               & 5000 & events recorded \\
\code{stop\_reason}         & reached 5000 events & why the run stopped \\
\code{dropped\_events}      & 0 & events lost by the recorder \\
\code{sampling\_mhz}        & 125.0 & sampling frequency \\
\code{tsample\_ns}          & 8.0 & sampling step, ns \\
\code{record\_length}       & 3000 & samples per waveform \\
\code{pretrigger}           & 200 & samples before the trigger \\
\code{adc\_nbit}            & 16 & ADC resolution \\
\code{adc\_to\_volts}       & $6.1455\times10^{-5}$ & volts per ADC count \\
\code{zin\_ohm}             & 100.0 & input impedance, $\Omega$ \\
\code{dc\_offsets\_percent} & [50.001, \dots] & DC offset per channel, \% of range \\
\code{vga\_gains\_db}       & [1, 1, 1, 1] & gain of the four input stages, dB \\
\code{baselines}            & [32777.6, \dots] & baseline per channel measured by WebDump with software triggers at the start of the run, ADC counts \\
\code{trigger\_mode}        & self & trigger mode \\
\code{trigger\_source}      & ITLA & trigger source (internal trigger logic A) \\
\code{trigger\_channels}    & "[0, 1, \dots, 63]" & channels in the trigger (a string) \\
\code{itl\_logic}           & Majority & logic of the internal trigger \\
\code{itl\_majority}        & 2 & channels required over threshold \\
\code{edge}                 & RISE & trigger edge \\
\code{threshold\_mode}      & Relative & thresholds relative to the measured baselines (or Absolute) \\
\code{thresholds}           & [19, \dots] & per-channel threshold as set, ADC counts \\
\code{thresholds\_board\_absolute} & [32796, \dots] & the same, as absolute ADC values loaded in the board \\
\bottomrule
\end{longtable}

\subsection{Conversions}

With $g$ = \code{adc\_to\_volts} and $Z$ = \code{zin\_ohm}, a sample of $a$
counts above baseline is a voltage $V = g\,a$, and the integral of a waveform
over a window, $\sum a_i\,\Delta t$ in ADC\,$\cdot$\,ns, is a charge
$Q = g\,\Delta t \sum a_i / Z$. For these runs $1$~ADC\,$\cdot$\,ns is
$6.1455\times10^{-4}$~pC. Pulses are positive (the trigger edge is
\code{RISE}).

\section{The LZF plugin}\label{sec:lzf}

WebDump compresses the waveforms with LZF, a filter that h5py carries built
in (HDF5 filter id 32000) but that is not part of the HDF5 library. The HDF5
library used by Julia (the \code{HDF5\_jll} artifact of \code{HDF5.jl},
version 2.2.0 at the time of writing) therefore cannot decompress the files
unless the filter is supplied as a plugin, a shared library that HDF5 loads
at run time from a directory in its plugin search path.

\code{src/h5lzf/} holds the six source files of the h5py filter
(\code{lzf\_filter.c}, \code{lzf\_filter.h} and the LZF library of Marc
Lehmann in \code{lzf/}, BSD licence in \code{LICENSE.txt}) and
\code{build.sh}, which compiles them into \code{libh5lzf.dylib} against the
headers and library of the HDF5 artifact that the Julia project uses. The
artifact is an MPI-enabled build, so its public header includes
\code{mpi.h}; the script finds that header in the matching MPI artifact of
the Julia depot. The plugin is specific to the HDF5 library it was built
against: rebuild it after the artifact changes. Python needs none of this.

\section{Module \code{JScan}}

Project at \code{src/jscan} (dependency: \code{HDF5.jl}). Load it with
\begin{lstlisting}
julia --project=src/jscan
julia> using JScan
julia> JScan.register_lzf()
julia> run = JScan.openrun("2026_10_02_bgo_cryo_tempscan/data_cryo/128K_37V_th19.h5")
julia> w = JScan.waveforms(run, 1:500)      # 3000 x 64 x 500 UInt16
julia> t = JScan.time_axis(run)             # ns from the trigger
julia> close(run.fid)
\end{lstlisting}

\subsection{\code{register\_lzf(dir = JScan.LZF\_PLUGIN\_DIR)}}

Adds \code{dir} to the plugin search path of the HDF5 library, unless it is
already there, and returns \code{dir}. The default is \code{src/h5lzf}. Throws
an \code{ArgumentError} if the directory does not exist or does not contain
\code{libh5lzf.dylib}. Must be called once per session before reading
waveforms; attributes and the small datasets can be read without it. The
plugin, once loaded, stays loaded for the session.

\subsection{\code{parse\_runname(name)}}

Parses a run name such as \code{"185K\_39d5V\_th19"} (no directory, no
extension) and returns the named tuple
\code{(T\_K = 185, V\_bias = 39.5, thr\_adc = 19)} with types
\code{Int}, \code{Float64}, \code{Int}. Throws an \code{ArgumentError} if the
name does not have the form \code{<int>K\_<int>[d<digits>]V\_th<int>}.

\subsection{\code{openrun(path)} and \code{openrun(f, path)}}

Opens the file read-only and returns the named tuple of
table~\ref{tab:run}. The file stays open, so that \code{waveforms} can read
from it; close it with \code{close(run.fid)}. The second form,
\code{openrun(path) do run ... end}, calls \code{f(run)} and closes the file
afterwards, also on error.

\begin{longtable}{l l p{78mm}}
\caption{Fields of the named tuple returned by \code{openrun}.}\label{tab:run}\\
\toprule
field & type & content, origin \\
\midrule
\endfirsthead
\toprule
field & type & content, origin \\
\midrule
\endhead
\code{name}   & \code{String} & run name, file name without extension \\
\code{path}   & \code{String} & path as given \\
\code{T\_K}   & \code{Int}     & crystal temperature, K; from the name \\
\code{V\_bias} & \code{Float64} & SiPM bias, V; from the name \\
\code{thr\_adc} & \code{Int}   & trigger threshold, ADC counts; from the name \\
\code{nev}    & \code{Int} & number of events; third dimension of \code{waveforms} \\
\code{nch}    & \code{Int} & number of recorded channels; second dimension \\
\code{nsamp}  & \code{Int} & samples per waveform; first dimension \\
\code{channels} & \code{Vector{Int}} & digitizer channel of each row; dataset \code{channels} \\
\code{tsample\_ns} & \code{Float64} & sampling step, ns; attribute \\
\code{sampling\_mhz} & \code{Float64} & sampling frequency; attribute \\
\code{pretrigger} & \code{Int} & pre-trigger samples; attribute. The trigger is at sample index \code{pretrigger + 1} (Julia, 1-based) \\
\code{adc\_to\_volts} & \code{Float64} & volts per ADC count; attribute \\
\code{zin\_ohm} & \code{Float64} & input impedance, $\Omega$; attribute \\
\code{baselines} & \code{Vector{Float64}} & WebDump baselines per channel, ADC counts; attribute \\
\code{thresholds} & \code{Vector{Int}} & per-channel thresholds as set; attribute \\
\code{threshold\_mode} & \code{String} & \code{"Relative"} or \code{"Absolute"}; attribute \\
\code{thresholds\_abs} & \code{Vector{Int}} & absolute thresholds in the board; attribute \code{thresholds\_board\_absolute} \\
\code{edge} & \code{String} & trigger edge; attribute \\
\code{itl\_logic} & \code{String} & internal trigger logic; attribute \\
\code{itl\_majority} & \code{Int} & majority level; attribute \\
\code{start\_time} & \code{String} & start of the run; attribute \\
\code{timestamp\_ns} & \code{Vector{Float64}} & trigger time of each event, ns; dataset \\
\code{trigger\_id} & \code{Vector{Int}} & board trigger counter per event; dataset \\
\code{fid} & \code{HDF5.File} & the open file \\
\code{wf} & \code{HDF5.Dataset} & the \code{waveforms} dataset, not read \\
\bottomrule
\end{longtable}

Attributes not in the tuple (board model, firmware, DC offsets, \dots) are
available as \code{HDF5.attrs(run.fid)}.

\subsection{\code{waveforms(run, events)}}

Reads the events in the range \code{events} (a \code{UnitRange}, 1-based)
and returns a (\code{nsamp}, \code{nch}, \code{length(events)}) array of
\code{UInt16}, raw ADC counts. With a single integer instead of a range, it
returns the (\code{nsamp}, \code{nch}) matrix of that event. Throws a
\code{BoundsError} if any event is outside \code{1:nev}. Reading is by whole
events, so blocks of a few hundred events are efficient (the file is chunked
by two events); a full run of 5000 events reads in about 1.3~s on an M-series
laptop.

\subsection{\code{time\_axis(run)}}

Returns a \code{Vector{Float64}} of length \code{nsamp} with the time of each
sample in ns from the trigger: $t_i = (i - 1 - \code{pretrigger})\,
\code{tsample\_ns}$, so \code{t[pretrigger + 1] == 0}.

\section{The pass}
\label{sec:pass}

\subsection{Purpose}

One pass over a run does everything that needs the raw per-channel
waveforms, so that they are never read again: the common-mode and baseline
correction and the summed waveform (\code{selection\_doc.tex}, steps 1 to
7), the charge and the time of the maximum of every channel in the
scintillation window (step~9, before any threshold), and the
single-photoelectron candidates of every channel in the pulse-free region
(\code{spe\_doc.tex}, steps 1 to 3). Everything else (classes, pile-up, zero
suppression, clusters, SPE spectra, equalization, light yield, decay time)
is computed offline in Python from the output, which is under 100~MB per
run. The definitions are those of the two documents cited; this section
gives the Julia conventions, the parameters, the output layout and the
tests.

\subsection{Usage}

\begin{lstlisting}
julia --project=src/jscan src/jscan/pass.jl RUN [options]
\end{lstlisting}

\code{RUN} is a run name or a path. Options, with the defaults of the
Python prototypes:

\begin{center}
\begin{tabular}{l l p{70mm}}
\toprule
option & default & meaning \\
\midrule
\code{--out DIR} & \code{pass/} (repository, not in git) & output directory; the file is \code{DIR/RUN\_pass.h5} \\
\code{--events a:b} & all & range of events, 1-based \\
\code{--template FILE} & \code{src/spe\_template\_RUN.csv} & SPE pulse shape (\code{template.py}); required \\
\code{--cm M} & \code{dead} & common-mode estimate: \code{dead}, \code{interp}, \code{median}, \code{none} (\code{selection\_doc.tex}, step~5) \\
\code{--dead LIST} & the twelve of \code{chanstat\_doc.tex} & channels whose mean is the common mode for \code{dead} \\
\code{--cm-smooth N} & 125 & moving-average width of that mean, in samples \\
\code{--edge E} & 32 & samples averaged at each end of $W$ for the interpolation \\
\code{--baseline-window N} & 250 & last $N$ samples: baseline of the sum \\
\code{--smooth N} & 64 & moving-average width of the pulse search \\
\code{--window A,B} & \code{1000,4000} & charge window, ns before and after the maximum \\
\code{--tail T} & 10 & SPE search region: last $T$ \us{} of the record \\
\code{--nsigma K} & 4.5 & SPE threshold in units of the filtered noise \\
\code{--block N} & 200 & events read per block \\
\bottomrule
\end{tabular}
\end{center}

\subsection{Per event}

Indices are 1-based; the trigger is at sample \code{pretrigger + 1}. For a
block of events the raw array is (samples, channels, events).

\begin{enumerate}
\item $x_{c,i} = w_{c,i} - b_c$ with $b_c$ the WebDump baseline (attribute).
\item $m_i = \operatorname{median}_c x_{c,i}$; provisional sum $S^0_i =
  \sum_c (x_{c,i} - m_i)$ minus its median over the last $N$ samples;
  moving average of width $N_{\rm sm}$ (edges use the samples available);
  $i^0_{\max}$ its maximum; $W = [i^0_{\max} - a, i^0_{\max} + b]$ clipped
  to the record.
\item $\tilde m$: with \code{dead}, the mean of $x_{c,i}$ over the dead
  channels, smoothed with a moving average of $N_{\rm cm}$ samples (no
  window involved); with \code{interp}, $m$ outside $W$ and a straight line
  inside between the means of $m$ over the $E$ samples before and after
  $W$ (one end replaces the other when $W$ touches the record edge); with
  \code{median}, $\tilde m = m$; with \code{none}, $\tilde m = 0$.
\item $y_{c,i} = x_{c,i} - \tilde m_i$; $S_i = \sum_c y_{c,i}$; $B$ =
  median of $S$ over the last $N$ samples; $S \leftarrow S - B$; $\sigma_1$
  = rms of $S$ there.
\item $\hat S$ = moving average of $S$; $i_{\max}$, $H = \hat S_{i_{\max}}$;
  $t_{\rm rise}$ = last sample before $i_{\max}$ with $\hat S < H/2$ (0 if
  none); charge window $[i_{\max} - a, i_{\max} + b]$ clipped, $n$ samples;
  control window as in \code{selection\_doc.tex} step~8.
\item Per channel: $B_c$ = median of $y_c$ over the last $N$ samples; $q_c
  = \Delta t \sum_{\rm window} (y_{c,i} - B_c)$; $q_c^{\rm ctrl}$ the same
  over the control window; $t_c$ = time of the maximum of the moving
  average (width $N_{\rm sm}$) of $y_c - B_c$ inside the window.
\item SPE, per channel, on $x_c$ minus its own median over the record: the
  matched filter $z_i = \sum_j T_j x_{c,i+j} / \sum_j T_j^2$ in the search
  region $R$ (last $T$ \us{}, $H$ samples trimmed at both ends, $H$ the
  template half length); $\sigma_z$ the clipped rms of $z$ over $R$
  (samples beyond three times a first rms removed); candidates the local
  maxima of $z$ above $K\sigma_z$, taken in decreasing order, at least
  $2H$ apart; for each, the time, $A = z_i$, and the residual rms of
  $x_{c,i+j} - A\,T_j$ over the window.
\end{enumerate}

\subsection{Output}

One HDF5 file per run, deflate level 4, shapes as seen from Julia (h5py
sees the axes reversed).

\begin{center}
\begin{tabular}{l l l p{62mm}}
\toprule
dataset & shape & type & content \\
\midrule
\code{events/index} & (nev,) & Int32 & event index, 0-based as in the run file \\
\code{events/trigger\_id}, \code{events/timestamp\_ns} & (nev,) & Int64, Float64 & copied from the run \\
\code{events/t\_max}, \code{events/H}, \code{events/t\_rise} & (nev,) & Float32 & ns from the trigger; height in ADC \\
\code{events/win\_lo}, \code{events/win\_hi} & (nev,) & Int32 & charge window, 0-based sample indices, inclusive \\
\code{events/ctrl\_lo}, \code{events/ctrl\_hi} & (nev,) & Int32 & control window \\
\code{events/baseline}, \code{events/rms} & (nev,) & Float32 & $B$ and $\sigma_1$ of the sum \\
\code{sum} & (nsamp, nev) & Int16 & $4\,S_i$ rounded (attribute \code{scale} = 4; values beyond $\pm 8191$ clipped and counted in the attribute \code{clipped}) \\
\code{qc}, \code{qc\_ctrl} & (nch, nev) & Float32 & $q_c$, $q_c^{\rm ctrl}$ in ADC\,ns, all channels, no threshold \\
\code{tc} & (nch, nev) & Float32 & $t_c$ in ns from the trigger \\
\code{spe/event}, \code{spe/channel} & (ncand,) & Int32, Int16 & one row per candidate \\
\code{spe/time}, \code{spe/A}, \code{spe/residual} & (ncand,) & Float32 & ns from the trigger, ADC, ADC \\
\code{spe\_sigma} & (nch, nev) & Float32 & $\sigma_z$ per channel and event \\
\bottomrule
\end{tabular}
\end{center}

Root attributes: the run name, path, T, bias, threshold, the run
attributes copied (\code{start\_time}, \code{pretrigger}, \code{tsample\_ns},
\code{baselines}), the thermometers T1, T2, T3 from \code{temperatures.txt}
when present, the true trigger rate from the trigger counter, every option
above, the template file and its values, the channel map file, the JScan
version, the git commit of the repository if available, and the date.
Equalization is not applied and dead channels are not masked: both are
offline decisions, and the per-channel charges are stored for all 64.

\subsection{Code}

\code{src/jscan/src/pass.jl}, included in the module: the pure functions
\code{common\_mode}, \code{interpolate\_cm}, \code{moving\_average},
\code{pulse\_search}, \code{windows}, \code{channel\_charges},
\code{matched\_filter}, \code{spe\_candidates}, one event's worth in
\code{process\_event}, and \code{pass(run, out; kwargs...)} which reads by
blocks and writes. \code{src/jscan/pass.jl} is the command-line script.

\subsection{Tests}

Two kinds. Unit tests of the pure functions on synthetic arrays: the
median of a known matrix; the interpolation with and without a record
edge; the moving average against a direct sum; a pulse placed at a known
sample found by \code{pulse\_search} with its rising edge; the windows at
the start, middle and end of a record; a template placed in zero noise
recovered by \code{matched\_filter} with amplitude one and residual zero;
candidates with the isolation rule.

An end-to-end test against the Python prototypes: a second fixture,
\code{test/data/real/197K\_39V\_th20.h5} (the run name, so that the reader parses it), holds three real events of
\code{197K\_39V\_th20} (events 0, 1 and 18: two REAL, one with pile-up),
cut by \code{test/make\_fixture\_real.py} with the full attributes of the
run; \code{selection.py --dump} and \code{spe.py --dump} write their
per-event quantities for the same file to JSON; the test runs \code{pass}
on the fixture and compares $t_{\max}$, $H$, $B$, $\sigma_1$, the window
bounds, $q_c$ for the 64 channels and the SPE candidates (count, times,
amplitudes) to the dump, to $10^{-3}$ relative for the floats and exactly
for the indices. Reading the real run \code{128K\_37V\_th19} for a few
events is kept as the integration test, skipped without the data symlink.

\section{Tests of the reader}

\code{src/jscan/test/runtests.jl}, run with \code{Pkg.test()}. They use a
fixture, \code{test/data/185K\_39d5V\_th19.h5}, with the layout of a real run
but 3 events, 64 channels and 300 samples, LZF-compressed, made once by
\code{test/make\_fixture.py} with h5py. The waveform values follow the
pattern $w[s, c, e] = 32768 + 100\,e + 3\,c + (s \bmod 11)$ with $e$, $c$,
$s$ zero-based, so reads can be checked exactly. Its attributes are those of
a real run, with \code{pretrigger} = 20 and \code{events} = 3.

\begin{itemize}
\item \code{parse\_runname}: integer bias (\code{37V}), half-volt bias
  (\code{39d5V}), thresholds 19 and 20; a malformed name throws.
\item \code{register\_lzf}: reading the fixture waveforms fails before
  registration (checked only when \code{HDF5\_PLUGIN\_PATH} is unset, since
  the plugin cannot be unloaded) and succeeds after; a second call is harmless;
  a missing directory throws.
\item \code{openrun}: every field of the tuple against the fixture's
  attributes and shape; the do-block form closes the file.
\item \code{waveforms}: shape, element type and exact values for an inner
  range, for the last block and for a single event; an out-of-range request
  throws.
\item \code{time\_axis}: length, zero at the trigger sample, step equal to
  \code{tsample\_ns}.
\item Integration: ten events of the real run \code{128K\_37V\_th19} with the
  expected dimensions; skipped if the data symlink is absent.
\end{itemize}

\end{document}
